\section{Conclusion and Future Work}
\label{sec:conclusion}

We have presented a production-sound integration of the WHIR
multilinear PCS into Ziren's STARK pipeline, closing all three
identified soundness gaps in the existing WHIR fast path.  Our
contributions span three layers:

\begin{itemize}
  \item \emph{Cryptographic protocols}: a sumcheck-based zerocheck
        for transition constraints, a sumcheck-corrected LogUp-GKR
        for lookup arguments, a late-binding WHIR adapter
        supporting post-commit evaluation points, a cross-binding
        spot-check protocol that ties FRI and WHIR commits to a
        single trace, and a Jagged + late-binding sumcheck reduction.

  \item \emph{Engineering integration}: opt-in env-flag wiring
        (\texttt{ZIREN\_USE\_WHIR=1},
        \texttt{ZIREN\_LATE\_BINDING=jagged},
        \texttt{ZIREN\_CROSS\_BINDING=1}) with default-preserved
        FRI behaviour; the recursion-circuit verifier scaffolding
        (covering all four protocol sub-steps); and a host-side
        prover/verifier path that integrates the new mechanisms
        without requiring upstream Plonky3 changes.

  \item \emph{Empirical validation}: end-to-end smoke tests on
        fibonacci, JSON, keccak, and large-sum (multi-shard)
        workloads, plus per-mechanism unit tests covering
        round-trips and tamper rejections.
\end{itemize}

\subsection{Future Work}

\paragraph{Recursion-circuit constraint emission --- orchestrator
shipped.}  \texttt{verify\_whir\_pcs\_in\_circuit} in
\texttt{crates/recursion/circuit/src/whir\_verifier.rs} composes
\texttt{emit\_sumcheck\_rounds}, \texttt{emit\_merkle\_path}, and
\texttt{emit\_final\_poly\_eval} into a single Builder<C>-emitting
orchestrator that mirrors
\texttt{p3\_whir::pcs::verifier::Verifier::verify} phase-by-phase:
(1)~observe initial commit + OOD answers,
(2)~initial sumcheck,
(3)~per WHIR round (tracking a rolling \texttt{prev\_commitment}):
observe new round commitment, observe OOD, sample STIR positions
via \texttt{challenger.sample\_bits}, verify each merkle path
against \emph{\texttt{prev\_commitment}} (not the new round
commitment), emit per-round sumcheck, roll \texttt{prev\_commitment
= round.commitment},
(4)~observe final-poly coefficients,
(5)~final STIR queries against the last \texttt{prev\_commitment},
(6)~assert \texttt{emit\_final\_poly\_eval} at the cumulative
challenges equals the running sumcheck claim, and
(7)~close each late-binding eq-statement claim.

A focused review identified and fixed one soundness bug in the
first draft: STIR queries had been asserted against
\texttt{round.commitment} (the NEW folded-poly root) instead of
\texttt{prev\_commitment} (the codeword actually being probed); a
malicious prover could have satisfied STIR by providing merkle
paths against a freshly-committed codeword while the openings came
from a divergent source.  Fixed by explicitly tracking
\texttt{prev\_commitment} and rolling it at the end of each round.

The orchestrator is intentionally simplified on two points relative
to upstream: no in-circuit \texttt{query\_pow\_witness} grinding
check (PoW bits are gated at the host), and the final
\emph{constraint-polynomial weight} check (\texttt{claimed\_eval
== weight * f(r)} in the upstream verifier) is approximated as a
direct final-eval equality, which is sound for our row-MLE usage
but not for general boundary constraints.  Extending to the full
\texttt{ConstraintPolyEvaluator} mirror is follow-up engineering
work.  Remaining integration steps: wire the orchestrator into the
recursion machine root verifier and feed a host \texttt{WhirProof}
through the witness stream as an \texttt{InCircuitWhirProof}.

\paragraph{Performance optimisation.}  Our profiling
(\S\ref{sec:performance}) shows that the WHIR open phase inside
\texttt{p3-whir}'s \texttt{Prover::prove} accounts for $97.9\%$
of jagged late-binding cost on a fibonacci shard.  The sumcheck
reduction we added is essentially free ($0.2\%$).  Upstream
parallelisation of \texttt{p3-whir}'s sumcheck rounds and merkle
folding is the highest-impact optimisation; alternatively,
splitting the jagged commit into multiple smaller WHIR proofs
trades a fraction of the proof-size win for wall-clock improvement.

\paragraph{Soundness analysis formalisation.}  The constructions
in this paper are presented operationally with informal soundness
sketches.  A formal soundness analysis (in the style of WHIR's
Theorem~5.2) for the late-binding adapter, the cross-binding
protocol, and the jagged sumcheck reduction would tighten the
security argument and enable third-party security audits.

\paragraph{Recursion-circuit \texttt{is\_div\_active} fix.}  The
preprocessed gate added to handle dead-DivF in \texttt{Select}
branches has a known interaction with \texttt{assert\_eq}'s
soundness DivF that bypasses the constraint when the witness
output cell is unread.  A focused fix --- either an
\texttt{is\_div\_soundness} flag or a Select-emission change ---
unblocks two pre-existing \texttt{should\_panic} unit tests and
removes a soundness footgun.

\paragraph{Production deployment.}  With the late-binding and
cross-binding chains complete on the host side, the next
deployment step is regenerating the gnark BN254 wrap circuit
artifacts (broken by the orthogonal Plonky3 FRI upgrade) and the
recursion compress \texttt{vk\_map.bin} (only required if the
recursion circuit changes propagate end-to-end).  The host-side
soundness work is independent of these regenerations.

\subsection{Acknowledgements}

This work builds on the open-source Plonky3 toolkit, the WHIR PCS
construction, and the LogUp-GKR lookup argument.  We thank the
authors of the SP1 hypercube architecture for the
\texttt{RecursiveWhirVerifier} reference implementation that
guided our recursion-circuit scaffolding port plan.
